\documentclass{article}
\usepackage{apacite}
\usepackage{hyperref}
\usepackage{multicol}
\usepackage{natbib}
\usepackage{sophia}
%\usepackage{elijah}
\usepackage[utf8]{inputenc}
\usepackage{array}
\setlength{\parindent}{0pt}
\usepackage{hyperref}

\usetikzlibrary{calc,patterns,angles,quotes}
\usepackage{pgfplots}
\pgfplotsset{width=15cm,height=10cm,compat=1.9}

\usepackage[english]{babel}
\usepackage{fancyhdr}
\usepackage{lastpage}
\pagestyle{fancy}
\fancyhf{}


\pagestyle{empty}% for cropping
\usepackage{tikz}
\usetikzlibrary{calc}
\usepackage{zref-savepos}
\usepackage{tabu}

\newcounter{NoTableEntry}
\renewcommand*{\theNoTableEntry}{NTE-\the\value{NoTableEntry}}

\newcommand*{\strike}[2]{%
  \multicolumn{1}{#1}{%
    \stepcounter{NoTableEntry}%
    \vadjust pre{\zsavepos{\theNoTableEntry t}}% top
    \vadjust{\zsavepos{\theNoTableEntry b}}% bottom
    \zsavepos{\theNoTableEntry l}% left
    \hspace{0pt plus 1filll}%
    #2% content
    \hspace{0pt plus 1filll}%
    \zsavepos{\theNoTableEntry r}% right
    \tikz[overlay]{%
      \draw
        let
          \n{llx}={\zposx{\theNoTableEntry l}sp-\zposx{\theNoTableEntry r}sp-\tabcolsep},
          \n{urx}={\tabcolsep},
          \n{lly}={\zposy{\theNoTableEntry b}sp-\zposy{\theNoTableEntry r}sp},
          \n{ury}={\zposy{\theNoTableEntry t}sp-\zposy{\theNoTableEntry r}sp}
        in
        (\n{llx}, \n{lly}) -- (\n{urx}, \n{ury})
      ;
    }% 
  }%
}

\usepackage{tikz}
\usetikzlibrary{calc}

\tikzset{
    right angle quadrant/.code={
        \pgfmathsetmacro\quadranta{{1,1,-1,-1}[#1-1]}     % Arrays for selecting quadrant
        \pgfmathsetmacro\quadrantb{{1,-1,-1,1}[#1-1]}},
    right angle quadrant=1, % Make sure it is set, even if not called explicitly
    right angle length/.code={\def\rightanglelength{#1}},   % Length of symbol
    right angle length=2ex, % Make sure it is set...
    right angle symbol/.style n args={3}{
        insert path={
            let \p0 = ($(#1)!(#3)!(#2)$) in     % Intersection**
                let \p1 = ($(\p0)!\quadranta*\rightanglelength!(#3)$), % Point on base line
                \p2 = ($(\p0)!\quadrantb*\rightanglelength!(#2)$) in % Point on perpendicular line
                let \p3 = ($(\p1)+(\p2)-(\p0)$) in  % Corner point of symbol
            (\p1) -- (\p3) -- (\p2)
        }
    }
}
\pagestyle{fancy}
\fancyhf{}
\title{Andi's Solution To BSS Math Team Selection Test}
\author{Andi Xie}
\date{October 6, 2022}
\lhead{Andi Xie}
\rhead{October 6, 2022}
\pagestyle{fancy}

 
\cfoot{Page \thepage \hspace{1pt} of \pageref{LastPage}}


\begin{document}

\maketitle


\section*{Question 1}
\begin{align*}
    a+c&=2b\\
    a=1,c&=3,5,7,9\\
    a=2,c&=4,6,8,10\\
    &\vdots
\end{align*}
\boxed{20} in total.


\section*{Question 2}
\begin{align*}
    \frac 43\pi r^3\left(75\%\right)&=\frac 13\left(\pi r^2\right)h\\
    \pi r^3&=\frac 13 \pi r^2h\\
    \frac hr&=\boxed{3}
\end{align*}



\section*{Question 3}
From Fermat's Little Theorem:
\begin{align*}
    1+29x^3-19x^2+199x+82x^2-75x+34x^3+1&\equiv 0\pmod{61}\\
    x^3+x^2+x+1&\equiv 0\pmod{61}\\
    (x+1)\left(x^2+1\right)&\equiv 0\pmod{61}\\
    x&=\boxed{11,50,60}
\end{align*}
Given that $x\in [0,60]$


\section*{Question 4}
\begin{align*}
    x^2&=\sqrt{13}x^{\log_{12}x}\\
    2\log_{13}x&=\frac 12\log_{13}13+\log_{13}x\left(\log_{13}x\right)\\
    2\left(\log_{13}x\right)^2-4\left(\log_{13}x\right)+1&=0\\
    \log_{13}x&=\frac{2\pm\sqrt 2}{2}\\
    \alpha\beta&=13^{\left(\frac{2-\sqrt 2}{2}\right)}\cdot 13^{\left(\frac{2+\sqrt 2}{2}\right)}\\
    \alpha\beta&=\boxed{169}
\end{align*}


\section*{Question 5}
Total - combinations with "JEFF"
Imagine "JEFF" as a block, permute with 4 other characters. That is $\frac{5!}{2!}-1$ combinations (1 from "JEFFJEFF"). \\
In total there are $\frac{8!}{2!2!4!}$ combinations, so
\[\frac{8!}{2!2!4!}-\frac{5!}{2!}+1=\boxed{361}\]


\pagebreak
\section*{Question 6}
Recall that the sum of the lengths of the two sides of a triangle has to be larger than the length of the third side. WLOG, let the length of the stick be 1, length of the segments be $x,y,1-x-y$. 

\begin{align*}
    x+y&>1-x-y\\
    y&<\frac 12-x\\
    x&<\frac 12\\
    y&<\frac 12
\end{align*}
Which is the desire cases. The total cases is $0<x+y<1$ which gives $y<-x+1$. Graph the equations:
\begin{center}
    \includegraphics[width=13cm]{Diagrams/stick.png}
\end{center}
We can see that the probability is $\boxed{\frac 14}$\\
\bigskip
\underline{Alternate Solution:}
Assume that $x+y>\frac 12$, which means that the second breaking point is pass the midpoint. 
\begin{align*}
    &2\int_{\frac{1}{2}}^1\left(x-\frac 12\right)dx\\
    =\ &2\left[\frac 12x^2-\frac 12x\right]_{\frac 12}^1\\
    =\ &\boxed{\frac 14}
\end{align*}


\section*{Question 7}
\begin{align*}
    F(x)&=\frac{\left(4x+\sqrt{4x^2-1}\right)\left(\sqrt{2x+1}-\sqrt{2x-1}\right)}{\left(\sqrt{2x+1}+\sqrt{2x-1}\right)\left(\sqrt{2x+1}-\sqrt{2x-1}\right)}\\
    &=\frac{4x\sqrt{2x+1}-4x\sqrt{2x-1}+\sqrt{8x^3+4x^2-2x-1}-\sqrt{8x^3-4x^2-2x+1}}{2}\\
    &=\frac{(2x+1)^{\frac 32}-(2x-1)^{\frac 32}}{2}\\
    \sum_{n=1}^{60}F(n)&=\frac 12\left(3^{\frac 32}-1^{\frac 32}+5^{\frac 32}-3^{\frac 32}+7^{\frac 32}-5^{\frac 32}+\cdots+119^{\frac 32}-117^{\frac 32}+121^{\frac 32}-119^{\frac 32}\right)\\
    &=\frac 12\left(121^{\frac 32}-1\right)\\
    &=\boxed{665}
\end{align*}


\section*{Question 8}
Each student sends out 15 slices and receives 15 slices. To avoid pizza-posers, student would receive pizza from those that he did not send to, which is 8 at most. 
\[\frac{(15-8) 24}{2}=\boxed{84}\]



\pagebreak
\section*{Question 9}
Consider the diamond shape $AE_1F_1BO$ (or you can think of it as $E_0E_1F_1F_0$), connect $E_1O$, $F_1O$, creating three isosceles triangles. $r=\frac{1}{\sqrt 2}$ from trivial Pythagorean. Maximizing the trapezoid is equivalent to maximize the diamond shape. Let $\theta$ to be the angle at the center of the triangles in sideways. The area of the diamond shape is given by:

\begin{align*}
    \left[AE_1F_1BO\right]&=r^2\sin\theta+\frac{r^2\sin\left(90^\circ-2\theta\right)}{2}\\
    &=r^2\sin\theta+\frac{r^2\left(1-2\sin^2\theta\right)}{2}\\
    &=r^2\sin\theta+\frac{r^2}{2}-r^2\sin^2\theta\\
    &=-\frac 12\left(\sin^2\theta-\sin\theta-\frac 12\right)
\end{align*}
To reach the maximum:
\begin{align*}
    \sin\theta&=\frac 12\\
    \theta&=30^\circ
\end{align*}
The angle is divided by 3 every time a trapezoid is made. Similarly, the next $\theta=10^\circ$. Generally, $\theta_n=\frac{90^\circ}{3^{n-1}}$. Therefore $\theta_{12}=\frac{10}{3^9}$. Divide by 3 gives $\angle E_{11}E_{12}O=\frac{10}{3^{10}}$. After some angle chasing, we find out that $N=\angle E_{12}E_{11}F_{11}=\frac{10}{3^{10}}$\\
$\therefore \frac 1N=\frac{3^{10}}{10}=\boxed{5904.9}$

\section*{Question 10}
Andi wrote nothing for Question 10.
\end{document}
